% =====================================================================================================
% Section I. The Training Problem (v3, after referee round 2)
% Writer: framework section writer (v3), 2026-09-24. Binding sources: paper/notes/revision_plan_v3.md (four main-text
% formal results: Lemma 1, Prop. 1 = generalized wedge incl. family budgets and the binding token cap; duality as text;
% notation q_T, q_D; "value of compactness", serving as one interpretation), paper/notes/theory_main_text_v2.md,
% appendix_proofs.tex (Lemmas A1-A5, Props. A8-A9, Cor. A3), round-2 reports R1-R4, memos rb1_sigmaC and rb2_decisions
% (+ reviews). v2 text: sections/v2/framework.tex.
% Numbers: reference technology = data/processed/ra2_wedge/headline.json (chin_q_point); SmolLM2 and Gopher wedges,
% sigma range, gamma: output/tables/fig1_geometry_numbers.csv (code/paper/make_fig1_geometry.py), which reproduces
% output/tables/ra2_wedge_models.csv (w_chin_q) to 1e-12. Figure 1 = output/figures/fig1_geometry.pdf.
% Number ledger, referee map and open issues: paper/notes/writer3_framework.md.
% Round 3 (WP3, fix list P1-P3, P10, P15): the value of compactness is the virtual value v_N (decision D-1); notation
% upsilon and varrho; Section I word count held at 2,300 by the offsetting cuts listed in paper/notes/round3_WP3_log.md.
% =====================================================================================================

\section{The Training Problem}\label{sec:framework}

This section writes a scaling law as a production technology with a known, multiplicative cost.

\subsection{Technology and Cost}\label{sec:framework:tech}

A developer trains a language model with $N$ parameters on $D$ tokens (tokens processed, counting repetitions) and
measures its cross-entropy loss $L$, in nats per token, on a held-out corpus. Following \citet{hoffmann2022training},
with productivity terms and an outer exponent added, we write
\begin{equation}\label{eq:tech}
L=E+e^{-\omega}\Big[A\big(e^{\psi_N}N\big)^{-\alpha}+B\big(e^{\psi_D}D\big)^{-\beta}\Big]^{\kappa}e^{\upsilon},
\qquad A,B,\alpha,\beta,\kappa>0 .
\end{equation}
Here $E$ is the irreducible loss, and the bracket collects the losses from finite parameters and data.
Productivity is Hicks-neutral ($\omega$) or factor-augmenting ($\psi_N$ and $\psi_D$; a cleaner corpus raises $\psi_D$),
and $\upsilon$ is seed and evaluation noise. The outer
exponent scales both output elasticities and leaves the isoquants unchanged, so it separates returns to scale from
substitution; $\kappa=1$ is testable (Section~\ref{sec:tech}).

Lower-case letters denote logs, $n\equiv\ln N$, $d\equiv\ln D$ and $c\equiv\ln C$, and $y\equiv-\ln(L-E)$ is the log
output index. With $u\equiv A(e^{\psi_N}N)^{-\alpha}$ and $v\equiv B(e^{\psi_D}D)^{-\beta}$, the output elasticities are
\begin{equation}\label{eq:elast}
\varepsilon_N\equiv\frac{\partial y}{\partial n}=\frac{\kappa\alpha u}{u+v},\qquad
\varepsilon_D\equiv\frac{\partial y}{\partial d}=\frac{\kappa\beta v}{u+v},
\end{equation}
and neither depends on $\omega$, $\upsilon$ or $E$. The family \eqref{eq:tech} is quasi-homothetic: the scaling
$(N,D)\mapsto(t^{\beta}N,t^{\alpha}D)$ multiplies reducible loss by $t^{-\kappa\alpha\beta}$ for every $t>0$ (Online
Appendix Lemma~\ref{lem:geometry}). Isoquants are therefore rescaled copies of one another, the compute-optimal mix
drifts with compute as a power law, and all curvature lies across the expansion path (Section~\ref{sec:ident}). This
restriction of the family, not a property of scaling laws, makes the curvature at the optimum the same at every compute,
whereas Section~\ref{sec:tech} finds $\sigma^*$ lower at the largest budgets of the two designs that reach them;
Lemma~\ref{lemma:sigma}, Proposition~\ref{prop:wedge} and the model-free results of Section~\ref{sec:ident} do not use
it.

Training a dense transformer takes about six floating-point operations (FLOPs) per parameter per token, so training
compute is $C\approx6ND$ \citep{kaplan2020scaling}. The approximation omits attention and embedding FLOPs,
and FLOPs are not dollars, so Proposition~\ref{prop:wedge} allows FLOP prices that vary with size and a cost of data.

In FLOP units the cost is multiplicative: both log-cost weights, $\partial c/\partial n$ and $\partial c/\partial d$,
equal one, isocosts are straight lines in logs, $n+d=c-\ln6$, and the implicit price of a parameter, $6D$, rises with the
quantity of the other input. Minimizing compute for a given loss therefore equates the output elasticities,
$\varepsilon_N=\varepsilon_D$, without the expenditure-share data that the first-order-condition approaches of
\citet{deloecker2012markups} and \citet{gandhi2020identification} require: FLOP accounting supplies the cost shares.

In this language, iso-loss curves are isoquants, an IsoFLOP profile traces loss along an isocost, the compute-optimal
allocation $\big(N^*(C),D^*(C)\big)$ is a conditional factor demand, and the loss--compute frontier $L^*(C)$ is the
inverse of a cost function (Figure~\ref{fig:geometry}, panel~a). We write $M\equiv D/N$ for tokens per parameter and
$M^*(C)$ for its compute-optimal value. Isoquants, the elasticity of substitution, the compute-optimal allocation and the
wedge defined below are ordinal: they are unchanged by any strictly increasing transformation of the same output (loss,
bits per byte, or accuracy on a benchmark that is monotone in loss) and do not depend on $E$, $\omega$ or $\kappa$
(Online Appendix Lemma~\ref{lem:sigma}). A capability that pretraining loss does not order, such as quality after
quantization \citep{kumar2024scaling}, is a different output with its own technology (Proposition~\ref{prop:wedge}(iii)). The frontier
elasticity is cardinal, and $M^*$ depends on the tokenizer and the parameter count. Online Appendix
Table~\ref{tab:app-dictionary} gives a full dictionary between scaling-law practice and production economics.

\subsection{Duality}\label{sec:framework:duality}

Let $a\equiv\beta/(\alpha+\beta)$, $\gamma\equiv\kappa\alpha\beta/(\alpha+\beta)$ and
$G\equiv(\alpha A/\beta B)^{1/(\alpha+\beta)}$, and set productivity and noise to zero. The allocation that minimizes
$6ND$ subject to $L\le\bar L$, or maximizes $y$ at given compute, satisfies $\varepsilon_N=\varepsilon_D$. The conditional
factor demands are
\begin{equation}\label{eq:demands}
N^*(C)=G\Big(\frac C6\Big)^{a},\qquad D^*(C)=G^{-1}\Big(\frac C6\Big)^{1-a},\qquad M^*(C)=G^{-2}\Big(\frac C6\Big)^{1-2a},
\end{equation}
and the loss--compute frontier and the cost function are
\begin{equation}\label{eq:frontier}
L^*(C)=E+K\Big(\frac C6\Big)^{-\gamma},\qquad C^*(L)=6\Big(\frac{K}{L-E}\Big)^{1/\gamma},
\qquad K\equiv\Big(\frac{\alpha+\beta}{\beta}\,AG^{-\alpha}\Big)^{\kappa}.
\end{equation}
Along the expansion path $\varepsilon_N=\varepsilon_D=\gamma$. At given compute, Hicks-neutral productivity leaves the
allocation unchanged, and factor-augmenting productivity moves it only through the bias index
$\chi\equiv\beta\psi_D-\alpha\psi_N$: $\ln M^*(C)$ falls by $2\chi/(\alpha+\beta)$ (Online Appendix
Lemma~\ref{lem:alloc}).

This is Shephard--Uzawa duality for the scaling law \citep{shephard1953cost,uzawa1964duality}: the technology, the
factor demands and the cost function carry the same information, linked by $a$ and $\gamma$. The outer exponent moves
the frontier but not the path. The frontier elasticity $\gamma$ is the returns-to-scale measure of
\citet{nerlove1963returns}; under the reference technology of Section~\ref{sec:wedge}, the $\kappa$-free fit to Chinchilla's 240 runs in
total parameters (Table~\ref{tab:designs}, panel~B), $\gamma=0.165$; because compute is quadratic in the scale of the
inputs, the technological scale elasticity on the path is $2\gamma$. A fit of \eqref{eq:tech} to all runs (Approach~3 of \citealp{hoffmann2022training}) is a
primal estimator; a regression of the loss-minimizing $\ln N$ of each IsoFLOP profile on $\ln C$ (their Approach~2)
estimates the factor demand, and a frontier through the same minima the cost function
(Online Appendix~\ref{app:additional} tests duality by comparing the two;
\citealp{klump2007factor} instead estimate a production function jointly with its first-order conditions). The minima are outcomes of inputs
that an experimenter assigned, not choices, so the endogeneity problems of IO arise for released models, not for
experiments.

\subsection{Substitution}\label{sec:framework:sigma}

The Hicks elasticity of substitution between parameters and data depends on the input mix only through the ratio of the
output elasticities, $w\equiv\varepsilon_N/\varepsilon_D=\alpha u/(\beta v)$, which we call the wedge. For every member
of the family \eqref{eq:tech},
\begin{equation}\label{eq:sigmaw}
\sigma(w)=\frac{1+w}{1+\alpha+w(1+\beta)},\qquad \sigma^*\equiv\sigma(1)=\frac{2}{2+\alpha+\beta},
\end{equation}
which follows from the formula of \citet{hao2026theory} (Online Appendix Lemma~\ref{lem:sigma}). The local elasticity
lies between $1/(1+\alpha)$ and $1/(1+\beta)$ and equals $\sigma^*$ on the compute-optimal path, where $w=1$. Under the
reference technology, $\sigma^*=0.701$, and $\sigma$ varies only between 0.699 and 0.702 over all input mixes.

\begin{lemma}[Interior compute optima require $0<\sigma<1$]\label{lemma:sigma}
Let $y=F(n,d)$ be any twice continuously differentiable technology with $F_n,F_d>0$, and let $C=6ND$. At an allocation
with $\varepsilon_N=\varepsilon_D$:
\begin{enumerate}
\item[(i)] the second-order condition for minimizing compute at given output holds strictly if and only if $0<\sigma<1$
there;
\item[(ii)] if $\sigma>1$ or $\sigma<0$ there, the allocation minimizes rather than maximizes output along its isocost;
if $\sigma>1$ everywhere, compute-optimal training is a corner solution at a bound on $N$ or $D$.
\end{enumerate}
Every member of the family \eqref{eq:tech} has $\sigma<1$ at every input mix, so the first-order condition
$\varepsilon_N=\varepsilon_D$ is sufficient.
\end{lemma}

\noindent The proof is in Online Appendix~\ref{app:proofs} (Lemma~\ref{lem:soc} and Corollary~\ref{cor:complements}).
With a linear cost $q_NN+q_DD$, isocosts are straight lines in levels, and any convex isoquant ($\sigma>0$) makes a
tangency a cost minimum. With $C=6ND$ they are straight lines in logs, and shifting toward one input lowers its own implicit price and
raises the other's, so cost alone favors a corner; an interior optimum requires the isoquant to be more curved than the
isocost, which is complementarity in this sense ($\sigma<1$).

The lemma's observable implication is that IsoFLOP profiles are U-shaped with an
interior minimum, as every design of Section~\ref{sec:tech} shows wherever its grid brackets the minimum. That
$\sigma^*<1$ is therefore not a finding; its level is. That $\sigma>0$ makes over-training a sensible response to a
value of compactness: in the Leontief limit of \citet{hao2026theory} the optimum sits at the kink whatever that value, whereas for $0<\sigma<1$ parameters can be traded for tokens at a cost governed by
$\sigma^*$ (Figure~\ref{fig:geometry}, panel~b).

\begin{figure}[tp]
\centering
\includegraphics[width=\textwidth]{fig1_geometry}
\caption{The Geometry of the Training Problem}\label{fig:geometry}
\begin{figurenotes}
Panel (a): the $(\log_{10}N,\log_{10}D)$ plane under the reference technology of Section~\ref{sec:wedge} (the $\kappa$-free fit
to Chinchilla's 240 runs, in total parameters; $\sigma^*=0.701$). Solid gray curves are isoquants at the frontier losses for
$C=10^{20},\dots,10^{24}$ FLOP; dotted lines are isocosts $6ND=C$; the blue line is the compute-optimal expansion path.
Filled orange dots are the three SmolLM2 base models (135M, 360M and 1.7B parameters, trained on 2, 4 and 11 trillion
tokens) at their total parameter counts, and hollow circles the same models at their non-embedding counts; for the
1.7B model the two markers nearly coincide, and the orange line starts at the total count. It joins SmolLM2 1.7B to
the path at equal compute; its wedge under the reference is $w=11.3$. Panel (b): three members
of the $\kappa$ family with $\sigma^*=0.60$, $0.74$ and $0.85$, constructed to share the expansion path and the
loss--compute frontier of the \citet{besiroglu2024chinchilla} parameters (Online Appendix Proposition~\ref{prop:fd}). The
vertical axis is reducible loss in excess of the frontier at $C=10^{22}$ FLOP, as a function of the deviation of tokens
per parameter from the compute-optimal ratio. The technologies coincide on the path and differ off it.
\end{figurenotes}
\end{figure}

\subsection{What Over-Training Reveals}\label{sec:framework:wedge}

Most open-weight language models released since 2023 lie far from the expansion path, with $M$ well above $M^*(C)$
(Section~\ref{sec:wedge}), whereas a developer that minimized training cost would have $w=1$.

\begin{proposition}[What over-training reveals]\label{prop:wedge}
A developer chooses $(N,D)$ to maximize $V\big(L(N,D),N\big)-X_T-q_DD$. Here $V$ is any differentiable objective that
falls with loss and may depend on size directly; $X_T\equiv q_T(N)\,6ND$ is training expenditure at the (shadow) price
$q_T(N)$ of a training FLOP, whose size elasticity is $\delta$; and $q_D\ge0$ is a cost per token. Let
$m_N\equiv-(\partial V/\partial\ln N)_L/X_T$ be the objective's own value of compactness: making the model one percent
smaller at the same loss is worth $m_N$ percent of training expenditure. Let $m_D\equiv q_DD/X_T$. Let $\nu_N$ be the
shadow value, per unit of training expenditure, of any binding constraint on size (a memory tier, a latency objective or
a size menu; $\nu_N\ge0$ for a cap, $\nu_N=0$ if none binds), and let $v_N\equiv m_N+\nu_N$ be the value of compactness:
the value at which a developer free of it would make the same choice, a virtual price in the sense of
\citet{neary1980theory}.
\begin{enumerate}
\item[(i)] (The wedge.) At an interior optimum with no binding cap on tokens,
\begin{equation}\label{eq:gw}
w=\frac{\varepsilon_N}{\varepsilon_D}=\frac{1+\delta+v_N}{1+m_D}.
\end{equation}
With $\delta=m_D=0$, $w-1=v_N$, and the compactness share $s\equiv(w-1)/w$ equals $v_N/(1+v_N)$.
\item[(ii)] (Serving.) Let $\delta=m_D=\nu_N=0$ and $V=R(L)-\lambda X_S$, so that serving is the only effect of size on $V$ at
given loss, where $X_S$ is the expenditure on serving $T$ tokens over the model's life at $2N$ FLOPs per token and a
constant $\varrho$ times the price of a training FLOP, with $T$ independent of size at given loss. Then
$w=1+\lambda\varrho T/(3D)$, $w-1=\lambda X_S/X_T$, and $s=\lambda X_S/(X_T+\lambda X_S)$ is the internalized serving
share of lifetime expenditure.
\item[(iii)] (Constraints and misspecification.) A binding cap on tokens adds $\mu\ge0$ to the denominator of
\eqref{eq:gw}, so with $\delta=m_D=0$, $w-1$ bounds $v_N$ from below; a binding cap on size enters $v_N$, so, if
$\mu=0$, $w-1$ bounds the objective's own $m_N$ from above. Logit distillation from a teacher with $N_T$ parameters is, if it leaves the
student's technology unchanged, a data cost with $m_D=N_T/(3N)$; if it also raises the productivity of tokens
($\chi>0$), it lowers the measured wedge further. If the developer's first-order conditions involve a technology $L'$
other than the econometrician's $L$ (its own productivity, a law it believes, or the output it values), the measured
wedge is $\hat w=w_{L'}\,w_L/w_{L'}$, where $w_{L'}$ is the wedge that the developer's technology assigns to its choice
and $w_L$ the one that the econometrician's assigns; it equals $w_{L'}e^{-\chi}$ under factor-augmenting productivity.
\item[(iv)] (Family token budgets.) Let a family train models of sizes $N_1,\dots,N_I$ on one token budget $D$, with
$\delta=m_D=0$. Member wedges then differ only through the technology; in family \eqref{eq:tech},
$\ln w_i-\ln w_j=-\alpha(\ln N_i-\ln N_j)$. Let $\bar w_H$ be the compute-weighted harmonic mean of the member wedges,
and let the weight $\pi_i$ be proportional to member $i$'s share of the family's training compute divided by $w_i$.
\begin{enumerate}
\item[(a)] If $D$ is chosen for the family, the family's first-order conditions identify one number,
$\sum_i\pi_iv_{N,i}=\bar w_H-1$, whether the sizes are chosen at the margin or set by binding size constraints.
\item[(b)] If a cap on tokens binds for the family, $\bar w_H-1\le\sum_i\pi_iv_{N,i}$.
\end{enumerate}
\end{enumerate}
\end{proposition}

\noindent The proof is in Online Appendix~\ref{app:proofs} (Propositions~\ref{prop:A-wedge} and~\ref{prop:A-family});
part (iv)(b) is Proposition~\ref{prop:A-family}(vi). The developer adds parameters
until their marginal value in loss equals their marginal cost: training expenditure plus whatever it loses when the model
is larger at the same quality (serving cost, adoption, latency, a memory tier). Tokens cost only training and data. An econometrician who sees only training cost therefore observes a
wedge: the developer acts as if parameters were $w$ times as expensive, relative to tokens, as training accounting says.

\textit{The value of compactness.} What $v_N$ measures depends on the developer's objective (Online Appendix
Table~\ref{tab:app-wedge-cases}). For a developer that serves this model, bears the cost ($\lambda=1$) and faces no
binding size constraint, $s$ is planned serving expenditure as a share of lifetime expenditure; a count of tokens,
$T=3D(w-1)/(\lambda\varrho)$, also needs the price ratio $\varrho$, likely above one because decoding is often bound by
memory bandwidth \citep{erdil2025inference}. For an open-weight developer whose users bear
the serving cost, $m_N$ is the value of compactness through adoption; under a binding tier, memory or latency cap
\citep{bian2025scaling}, the value of compactness includes the cap's shadow value. Section~\ref{sec:wedge} tests these
readings as a problem of conduct \citep{berry2014identification,backus2021common,duarte2024testing}, and the evidence is
more consistent with adoption and memory tiers than with the internalization of serving cost. We therefore call $v_N$
the value of compactness and $s$ the compactness share throughout, and read $s$ as planned serving expenditure only
where the developer served the size and neither a common token budget nor a tier window intervenes
(Section~\ref{sec:wedge-conduct}).

\textit{Other sources of a wedge.} Their signs are known (Online Appendix Proposition~\ref{prop:A-wedge}(iv)). A price
of a FLOP that rises with size ($\delta>0$) and parameter-augmenting productivity ($\chi<0$) raise the measured wedge.
So does post-training compute that scales with $N$ and is planned when $N$ is chosen: counting its expenditure $X_P$ as a
cost alongside $X_T$, with $x_P\equiv X_P/X_T$, $w=1+x_P+m_N$, where $m_N$ excludes post-training, and
Section~\ref{sec:wedge} bounds $x_P$ with disclosed budgets. Data costs, caps on tokens, distillation
\citep{busbridge2025distillation} and data-augmenting productivity ($\chi>0$) lower the wedge. Serving demand cannot
produce a wedge below one as long as serving cost rises with size. Gopher's $w=0.28$ under the reference technology
\citep{rae2021scaling} is thus consistent with a binding data constraint or a data cost, better data, a size target (a
minimum-size floor), or optimization against another technology, plausibly the law of \citet{kaplan2020scaling}.

\textit{Family budgets.} On one token budget, member wedges are fixed by the technology and are not revealed
preferences. A binding cap, as when every size trains on the whole curated corpus, makes the family share $1-1/\bar w_H$ a lower bound,
and a member's own wedge one only if that member would have trained on more tokens at its own size, which a cap on the
family's budget does not imply.

\textit{A closed form.} In the $\kappa$ family the wedge of a model with compute $C$ and tokens per parameter $M$ is
\begin{equation}\label{eq:wM}
\ln w=\Big(\frac1{\sigma^*}-1\Big)\ln\frac{M}{M^*(C)},
\end{equation}
where $M^*(C)$ is the compute-optimal ratio at the model's own compute (Online Appendix Lemma~\ref{lem:geometry} and
Corollary~\ref{cor:suff}). The log wedge is the distance from the expansion path at equal compute, scaled by curvature
(Figure~\ref{fig:geometry}, panel~a); the technology supplies the zero point $M^*(C)$ and the scale $1/\sigma^*-1$.
SmolLM2 1.7B \citep{benallal2025smollm2}, trained on 11 trillion tokens, has
$M\approx6{,}400$, about 290 times the reference technology's compute-optimal ratio at its compute. With the reference
scale, $1/\sigma^*-1=0.427$, its wedge is $w\approx11.3$: its developer behaves as if making it one percent smaller at
the same loss were worth about 10 percent of the training bill, and $s\approx0.91$. Such numbers are only as good as
the technology's two inputs, neither of which any public design observes at SmolLM2's compute. The sign needs only the
zero point: for any technology with single-peaked IsoFLOP profiles, $w>1$ if and only if $M>M^*(C)$, whatever the
scale, so bounds on $M^*(C)$ identify the sign (Proposition~\ref{prop:pi}).
